\section{Introduction}

A pension that is withdrawn as the recipient's assets rise gives the household a reason to hold
fewer assets, and it is often said that this makes the household's policy function non-monotone.
The claim is usually supported by an appeal to the value function: the means test puts a kink, or
a cliff, into the reward from saving, the value function is no longer concave, and monotone
comparative statics are said to fail with it.

Three quotations fix what is actually claimed. \citet{KudrnaWoodland2011}, simulating the
Australian age pension in a general equilibrium life-cycle model, treat the geometry as settled:
``the means test for the age pension causes the budget set to be non-convex, as is well known.''
They pay for it in their solution method, whose solver ``is not guaranteed to yield a globally
optimal solution for the inter-temporal expected utility maximization problem when the budget
constraint is non-convex''; they check the solution of every household over the age of sixty by
hand to confirm that the reported optimum is global. \citet{GarstenauerSiassi2024}, modelling the
United States tax and transfer system for married couples with children, reach the same conclusion
about method: ``the presence of kinks and non-differentiabilities in the budget constraint
generated by the tax-transfer system renders Euler equation methods inapplicable.'' They solve by
discrete-state value function iteration, and because ``of the presence of these kinks and
non-convexities, it is crucial for us to capture as accurately as possible how labor supply
choices respond to taxes and transfers'' they put mothers' annual hours on a grid whose adjacent
nodes are twenty-five hours apart.

The third is closest to our question, and it cuts the other way. \citet{Wellschmied2021} studies
an asset means test in a life-cycle model and characterises the policy analytically. His Theorem 1
states that the saving policy ``is increasing'', and his proof that the value function stays
sub-differentiable uses that the policy ``is weakly increasing''. What he shows fails is
everything else. ``Despite the policy $k_h(\cdot,X)$ being monotone, the marginal propensity to
consume out of additional assets may be one for households with low income and assets''; ``the
policy function makes a jump''; and ``choices are distorted, despite satisfying first order
conditions because of the kinks in the expected value function''. A footnote records how fast this
compounds, since there is ``a rapid increase in the number of non-differentiabilities in the value
function''. The cost lands in the numerical appendix, where ``first order conditions are not
necessarily unique'', so off-grid choice is abandoned over the part of the state space the test
touches and almost nine thousand asset nodes are carried instead. His argument works through the
differential structure of the value function, following \citet{ClausenStrub2020} on how kinks
propagate under maximisation; the argument below uses no calculus at all.

The conclusion is right and the reason is wrong, and \citet{Wellschmied2021} is already halfway to
saying so. Monotonicity of the saving policy in cash on hand
survives \emph{any} continuation value whatever --- concave, kinked, discontinuous, or produced by
a means test of any shape --- because the only structure the comparison needs is strict concavity
of period utility, and the continuation cancels between the two optimality inequalities. This is
Topkis's argument \citep{Topkis1978,MilgromShannon1994} applied to the household's problem, and it
does not care what the continuation looks like.

What a non-concave continuation does produce is a jump. At the wealth where the household stops
protecting its pension and starts saving past the test, optimal saving jumps up, and because it
jumps by more than wealth has risen, consumption jumps \emph{down}. So the non-monotonicity is
real, it is in consumption, and it coexists with a perfectly monotone saving policy. We give an
exact two-period example, with the optimum at each of two wealth levels computed and verified.

Which policy misbehaves also depends on which state space one works in, and the two are not
equivalent under a means test. In the household's own state space --- assets $a$ and the earnings
shock $z$ --- next period's assets rise with this period's at each $z$, again for every
continuation. The Lean development that this paper draws on already had that statement, but with
concavity of the continuation as a hypothesis; a means test is exactly what removes that
hypothesis's premise, so we prove it without. Monotone saving in assets then fails in one case
only: when the test is assessed on the assets currently held and withdraws the pension faster than
the gross return, cash on hand itself falls as assets rise, and the policy follows it down. The
threshold is sharp, the taper rate against the gross return, and Australia's asset test sits on
the benign side of it.

Section~\ref{sec:model} sets out the retired household and the three pension schedules.
Section~\ref{sec:cash} gives the two results in the cash-on-hand state space: saving monotone
always, consumption not. Section~\ref{sec:assets} gives the result in the $(a,z)$ state space and
the exception. Section~\ref{sec:num} calibrates a retiree and checks each case. Every result
stated with a formal name has been checked by the Lean 4 proof assistant \citep{Lean4} against
Mathlib \citep{Mathlib}; the development is at
\url{https://github.com/LeanEconomics/LeanEconomics}, files
\texttt{LeanEconomics/OLG/MeansTest.lean} and \texttt{Retirement.lean}.
